\chapter{Teaching A Machine}
\label{stone:computer}

By now we have built a complete arithmetic. We can measure. We can join constructions. We can let them grow. We can multiply. We can divide. We can describe space. We can reflect it. Everything has been done with exact arithmetic.

Now a different question appears. Can a machine do the same thing?

\section*{Machines Do Not Think}

People often speak about computers as though they think. They do not. A computer follows instructions. Exactly. Patiently. Relentlessly. If the instructions are correct, the computer produces the correct result. If the instructions are wrong, it faithfully produces the wrong result.

A computer has no opinion. It has no intuition. It has no understanding. It simply performs the arithmetic we give it. That means the quality of the mathematics matters more than the speed of the processor. We do the mathematics. The computer does the work.

\section*{Arithmetic Is Work}

Every computer program eventually becomes a long sequence of tiny operations. Add. Subtract. Multiply. Divide. Compare. Repeat. Millions of times. Sometimes billions. The remarkable things computers accomplish are built from these tiny steps. If we improve the arithmetic, every larger computation improves with it.

\section*{A Different Philosophy}

Most attempts to make computers faster begin with the machine. Build a faster processor. Build a wider memory. Increase the clock speed. PhiBase begins somewhere else. Can we give the processor less work to do?

That is a builder's question. Builders rarely become more productive by working harder. They become more productive by working smarter.

\section*{Remember What Already Works}

Earlier we built the Fibonacci Spine. Then the Dense Grid. Then the Digital Slide Rule. All three share the same idea. Never rediscover something you already know. A computer can follow exactly the same philosophy. If a result has already been computed, remember it. Retrieve it. Continue building. The arithmetic becomes simpler because the difficult work has already been done.

\section*{Multiplication Without Multiplying}

One of the oldest microprocessors ever built, the MOS 6502, had no hardware multiplication instruction. It survived for decades because engineers discovered that many multiplications could be replaced by carefully prepared lookup tables. The work moved from computation into memory.

PhiBase follows exactly the same engineering instinct. Whenever practical, the difficult arithmetic is performed once. Later computations simply reuse that knowledge. The mathematics stays exact. The machine performs less work.

\section*{The Experiment}

Eventually every idea reaches the same moment. Someone builds it. A complete transformer inference engine was rewritten to use PhiBase arithmetic. The model itself did not change. The arithmetic did. The same questions were asked. The same answers were produced. The processor simply performed different arithmetic while doing its work.

The result was encouraging. The outputs closely matched those of the original implementation while requiring substantially fewer arithmetic cycles. The experiment did not prove that every future computer should use PhiBase. It answered a simpler question. Can a modern machine work this way? Yes.

\section*{Why This Matters}

Artificial intelligence is only one example. Robotic construction is another. Architectural simulation is another. Crystal growth. Structural analysis. Any computation built upon exact golden-ratio geometry becomes a candidate. PhiBase is not trying to replace every arithmetic. It is offering a better arithmetic wherever this geometry naturally appears.

\section*{Builder's Notebook}

Imagine two workshops. The first measures every board every morning. The second keeps carefully labeled templates for every common part. Both produce the same furniture. Which one would you rather work in? Why? Now replace the boards with arithmetic. The question remains the same.

\section*{Looking Ahead}

This book began with two triangles. It ended with a machine capable of using the arithmetic those triangles inspired. That may seem like a very long journey. It is. But the most surprising discovery was not the arithmetic. It was the way the arithmetic was discovered. That is our final stepping stone.
